%Tiefe des Inhaltsverzeichnisses setzen:
% bei Report: 3 geht bis subsection
\newcommand{\mytocdepth}{2}
%Tiefe des Inhaltsverzeichnisses in PDF-Lesezeichen setzen (evtl. höher als die des toc)
%\newcommand{\mybookmarkdepth}{\mytocdepth}
\newcommand{\mybookmarkdepth}{3}
%
% Hier Breite von marginnotes setzen
\newlength{\mymarginparwidth}
%\setlength{\mymarginparwidth}{0.15\paperwidth}
\setlength{\mymarginparwidth}{23.1mm}%was 10mm
%\setlength{\mymarginparwidth}{0pt}
% Abstand zwischen Seitennotizen und Text setzen
\newlength{\mymarginparsep}
%\setlength{\mymarginparsep}{}
\setlength{\mymarginparsep}{1mm}%was 3mm
%\setlength{\mymarginparsep}{5mm}
%\setlength{\mymarginparsep}{0pt}




%%% Absatzangaben von Komaskript:
%%%\setparsizes{<Einzug>}{<Abstand>}{<Endzeilenleerraum>}
%%%\setparsizes{<parindent>}{<parskip>}{<parfillskip>}
%%\setparsizes{\myparindent}{\myparskip}{0pt plus\dimexpr\textwidth-10ex}
%\setparsizes{0pt}{6.6pt}{0pt plus\dimexpr\textwidth-10ex}




% Abstand zwischen den Paragraphen
\newlength{\myparskip}
%\setlength{\myparskip}{0.5\baselineskip plus 0.5\baselineskip} % das wäre wohl default
%\setlength{\myparskip}{0.5\baselineskip} %ohne Pluswerte (= keine Dehnung) und Minuswerte (= keine Stauchung)
%\setlength{\myparskip}{0.6\baselineskip} %ohne Pluswerte (= keine Dehnung) und Minuswerte (= keine Stauchung)
%\setlength{\myparskip}{0.6\baselineskip plus 1pt minus 1pt} %mit Pluswerten (= Dehnung) und Minuswerten (= Stauchung)
%\setlength{\myparskip}{0.6\baselineskip plus 0pt minus 0pt} %mit Pluswerten (= Dehnung) und Minuswerten (= Stauchung)
\setlength{\myparskip}{6.6pt} %ohne Pluswerte (= Dehnung) und Minuswerte (= Stauchung)

\newlength{\mysmallskip} % default: 3.0pt plus 1.0pt minus 1.0pt
\setlength{\mysmallskip}{3pt} %ohne Pluswerte (= Dehnung) und Minuswerte (= Stauchung)

\newlength{\mymedskip} % default 6.0pt plus 2.0pt minus 2.0pt
%\setlength{\mymedskip}{6pt plus 1pt minus 1pt} %mit Pluswerten (= Dehnung) und Minuswerten (= Stauchung)
\setlength{\mymedskip}{6pt} %ohne Pluswerte (= Dehnung) und Minuswerte (= Stauchung)

\newlength{\mybigskip} % default 12.0pt plus 4.0pt minus 4.0pt
%\setlength{\mybigskip}{12pt plus 2pt minus 2pt} %mit Pluswerten (= Dehnung) und Minuswerten (= Stauchung)
\setlength{\mybigskip}{12pt} %ohne Pluswerte (= Dehnung) und Minuswerte (= Stauchung)

%%% Vertikaler Abstand zwischen zwei Gleitobjekten
\newlength{\myfloatsep}
%\setlength{\myfloatsep}{\baselineskip} % (default 12pt plus 2pt minus 2pt)
%\setlength{\myfloatsep}{\baselineskip}
\setlength{\myfloatsep}{12pt plus 2pt minus 2pt}

%%% Vertikaler Abstand zwischen der Top- oder Bottom-Bereich und Text-Bereich
\newlength{\mytextfloatsep}
%\setlength{\mytextfloatsep}{\baselineskip} % (default 20pt plus 2pt minus 4pt)
%\setlength{\mytextfloatsep}{\myfloatsep}
%\setlength{\mytextfloatsep}{1.5\baselineskip}
\setlength{\mytextfloatsep}{15pt plus 2pt minus 4pt}

%% Abstand zwischen Inline-Gleitobjekten, die mit "here" platziert worden sind,
%% und dem darüber und darunter angeordneten Fließtext fest
%% Beispiel: \intextsep5mm plus3mm minus2mm
%\setlength{\intextsep}{0.5\baselineskip} % Platz ober- und unterhalb des Bildes
\newlength{\myintextsep}
%\setlength{\myintextsep}{2\baselineskip} % (default 12pt plus 2pt minus 2pt) 
\setlength{\myintextsep}{15pt plus 2pt minus 2pt} % KSP, real ca. 7 mm vor und nach Listings
%\setlength{\myintextsep}{10pt}
%\setlength{\myintextsep}{12pt} %ohne Pluswerte (= Dehnung) und Minuswerte (= Stauchung)
%\setlength{\myintextsep}{12pt plus 1pt minus 1pt} %mit Pluswerten (= Dehnung) und Minuswerten (= Stauchung)
%\setlength{\myintextsep}{0pt} %kein Zusatzabstand!


\newlength{\mychapterbeforeskip}
\newlength{\mychapterafterskip}
\newlength{\mysectionbeforeskip}
\newlength{\mysectionafterskip}
\newlength{\mysubsectionbeforeskip}
\newlength{\mysubsectionafterskip}
\newlength{\mysubsubsectionbeforeskip}
\newlength{\mysubsubsectionafterskip}
\newlength{\myparagraphbeforeskip}
\newlength{\myparagraphafterskip}

\setlength{\mychapterbeforeskip}{8\myparskip}
\setlength{\mychapterafterskip}{2.5\myparskip}
%
\setlength{\mysectionbeforeskip}{2\myparskip}
\setlength{\mysectionafterskip}{1.5\myparskip}
%
\setlength{\mysubsectionbeforeskip}{1.5\myparskip}
\setlength{\mysubsectionafterskip}{1\myparskip}
%
\setlength{\mysubsubsectionbeforeskip}{1.5\myparskip}
\setlength{\mysubsubsectionafterskip}{1\myparskip}
%
\setlength{\myparagraphbeforeskip}{1.5\myparskip}
\setlength{\myparagraphafterskip}{1\myparskip}


\newlength{\mylistingframesep}
\setlength{\mylistingframesep}{1ex}

\newlength{\mylistingframerule}
\setlength{\mylistingframerule}{0.4pt}


%%Abstand zwischen der Abbilung und der Überschrift:
%%funktioniert aus irgendeinem Grund nicht!
%%--> reason is "\DeclareFloatVCode{negintextskip}{{\vspace{-\myintextsep}\par}}" and its use in for figures and tables
\newlength{\myaboveskip}
%\setlength{\myaboveskip}{10pt}
\setlength{\myaboveskip}{0pt}

%%Abstand nach der Überschrift einer Abbilung:
%%funktioniert aus irgendeinem Grund nicht!
\newlength{\mybelowskip}
%\setlength{\mybelowskip}{10pt}
\setlength{\mybelowskip}{0pt}


%Abstand zwischen der Subfigure-Überschrift und der Überschrift:
\newlength{\mysubcaptionskip}
\setlength{\mysubcaptionskip}{6pt}

%Abstand zwischen der Float-Umgebung und der Unter-/Überschrift
\newlength{\mycaptionskip}
\setlength{\mycaptionskip}{10pt}


%Abstand zwischen der Abbilung/Tabelle und der Überschrift:
\newlength{\myfigureskip}
\setlength{\myfigureskip}{\mycaptionskip}
%Abstand zwischen der Abbilung und der Überschrift:
%funktioniert aus irgendeinem Grund nicht!
\newlength{\myfigureaboveskip}
\setlength{\myfigureaboveskip}{\myaboveskip}
%Abstand nach der Überschrift einer Abbilung:
\newlength{\myfigurebelowskip}
\setlength{\myfigurebelowskip}{\mybelowskip}
%Abstand zwischen der Abbilung und der Überschrift:
\newlength{\myfigurecaptionskip}
\setlength{\myfigurecaptionskip}{\mycaptionskip}



%Abstand zwischen der Abbilung/Tabelle und der Überschrift:
\newlength{\mytableskip}
\setlength{\mytableskip}{\mycaptionskip}
%Abstand zwischen der Tabelle und der Überschrift:
%funktioniert aus irgendeinem Grund nicht!
\newlength{\mytableaboveskip}
\setlength{\mytableaboveskip}{\myaboveskip}
%Abstand nach der Überschrift einer Tabelle:
\newlength{\mytablebelowskip}
\setlength{\mytablebelowskip}{\mybelowskip}
%Abstand zwischen der Tabelle und der Überschrift:
\newlength{\mytablecaptionskip}
\setlength{\mytablecaptionskip}{\mycaptionskip}


%% Abstand zwischen der Listing und der Überschrift:
\newlength{\mylistingskip}
\setlength{\mylistingskip}{\mycaptionskip}
%% Abstand zwischen dem Text und der Listing über der Listing (d.h. ggf. über der Überschrift)
\newlength{\mylistingaboveskip}
%\setlength{\mylistingaboveskip}{\myaboveskip}
%\setlength{\mylistingaboveskip}{0.5\mytextfloatsep}
%\setlength{\mylistingaboveskip}{0pt}
\setlength{\mylistingaboveskip}{\myintextsep}
%% Abstand nach der Listing (d.h. ggf. unter der Unterschrift):
\newlength{\mylistingbelowskip}
%\setlength{\mylistingbelowskip}{\mybelowskip}
%\setlength{\mylistingbelowskip}{\mytextfloatsep}
%\setlength{\mylistingbelowskip}{0pt}
\setlength{\mylistingbelowskip}{\myintextsep}
%\setlength{\mylistingbelowskip}{\myintextsep}
%\setlength{\mylistingbelowskip}{2\parskip}
%%% Abstand zwischen der Listing und der Überschrift:
%\newlength{\mylistingcaptionskip}
%\setlength{\mylistingcaptionskip}{\mycaptionskip}
%%% Abstand vor der Überschrift einer Listing:
%\newlength{\mylistingcaptionaboveskip}
%\setlength{\mylistingcaptionaboveskip}{\myparskip}
%%% Abstand nach der Überschrift einer Listing:
%\newlength{\mylistingcaptionbelowskip}
%\setlength{\mylistingcaptionbelowskip}{\myparskip}



%Abstand zwischen zwei Subfloats (muss manuell hinzugefügt werden)
\newlength{\mysubfloatvskip}
\setlength{\mysubfloatvskip}{10pt} % (default 12pt plus 2pt minus 2pt)


%Abstand oberhalb von Subfigures:
\newlength{\mysubfigaboveskip}
%% Problem: Fügt auch einen (unnötigen) Abstand vor der ersten Reihe von Subfigures ein,
%% so dass der Abstand zwischen dem vorherigen Absatz und der Abbildung unnötig vergrößert wird.
%% Dasselbe gilt auch am Anfang einer Seite, auch wenn die Doku sagt, dass dort der Abstand unterdrückt werden sollte 
%% Muss daher auf Null gesetzt werden
%\setlength{\mysubfigaboveskip}{\mysubfloatvskip} %funktioniert nicht
\setlength{\mysubfigaboveskip}{0pt}
%Abstand nach der Subfigure-Unterschrift
\newlength{\mysubfigbelowskip}
%% Problem: Fügt auch einen (unnötigen) Abstand nach der letzten Reihe von Subfigures ein,
%% so dass der Abstand zwischen den Subfigure-Unterschriften der letzten Subfigure-Reihe
%% und der Gesamt-Unterschrift der Abbildung zu groß wird!
%% Muss daher auf Null gesetzt werden
%\setlength{\mysubfigbelowskip}{\mysubfloatvskip} %funktioniert nicht
\setlength{\mysubfigbelowskip}{0pt}
%Abstand zwischen der Abbilung/Tabelle und der Überschrift:
\newlength{\mysubfigcaptionskip}
\setlength{\mysubfigcaptionskip}{\mysubcaptionskip}




%% Weitere Angaben:
%% Forderung des KSP-Verlages:
%usepackage[a5paper,headheight=1.5\baselineskip,top=25mm,lines=31,heightrounded=true,bindingoffset=15mm,textwidth=106mm]{geometry}
%\usepackage[a4paper,headheight=1.5\baselineskip,top=25mm,lines=46,heightrounded=true,bindingoffset=15mm,textwidth=160mm]{geometry}
% footskip: Abstand zwischen Textunterkante und Seitenzahl
% KIT-Verlag: mindestens 10mm bzw. 3 Zeilen
% Koma-Standard 3.5\baselineskip
% IES-Vorlage: 2\baselineskip
% Abstand zwischen Textkörper und Unterkante Fußzeile (Seitenzahlen)
\newlength{\myfootskip}
\setlength{\myfootskip}{11mm}
% Abstand zwischen den Fußnoten und dem Text
\newlength{\myfootnoteskip}
\setlength{\myfootnoteskip}{1.5\baselineskip}
% Abstand zwischen Textkörper und Linie in der Kopfzeile
\newlength{\myheadsep}
%\setlength{\myheadsep}{5mm}
\setlength{\myheadsep}{7mm} %Da eine kleinere Schriftgröße verwendet wird

\newlength{\myheadheight}
\setlength{\myheadheight}{1.5\baselineskip}

\newlength{\mytop}
\setlength{\mytop}{25mm}

\newlength{\mybindingoffset}
\setlength{\mybindingoffset}{0mm}

\newlength{\myinner}
\setlength{\myinner}{25mm}%a4 with 160mm textwidth has 50mm for inner+outer length

\newlength{\myouter}
\setlength{\myouter}{25mm}%a4 with 160mm textwidth has 50mm for inner+outer length

\newlength{\mytextwidth}
\setlength{\mytextwidth}{160mm}%160mm for a4, 106mm for a5

\newlength{\mypagewidth}
\newlength{\mypageheight}
\setlength{\mypagewidth}{\paperwidth}
\setlength{\mypageheight}{\paperheight}

\newlength{\effectivecenterx}
\setlength{\effectivecenterx}{0.5\mypagewidth + 0.5\mybindingoffset}

\newlength{\effectivecentery}
\setlength{\effectivecentery}{0.5\mypageheight}

%vertical float alignment for the only-float pages:
%on float-only pages the figures/tables shall be aligned at the top of the page instead of being centered
\newboolean{SetFloatsVerticallyCentered}
\setboolean{SetFloatsVerticallyCentered}{false}